\documentclass[10pt,a4paper]{amsart}
\usepackage[margin=19mm]{geometry}
\usepackage{amsmath,amssymb,mathtools}
\usepackage[scr=rsfs,cal=euler]{mathalfa}
\usepackage{xfrac}
\usepackage[unicode,hidelinks,bookmarks=false]{hyperref}
\newcommand{\ie}{i.e.\ }
\newcommand{\cf}{cf.\ }
\newcommand{\resp}{respectively\ }
\newcommand{\Subsection}{\subsection}
\providecommand{\nobreakdash}{\nobreak\hbox{-}}
\setlength{\emergencystretch}{2em}
\multlinegap0pt
\usepackage{bm}
\usepackage{bbm}
\usepackage{dsfont}
\usepackage{xspace}
\usepackage{cancel}
\usepackage{mathtools}
\usepackage{cleveref}
\usepackage[shortlabels]{enumitem}
\usepackage{amsthm}
\makeatletter
\@ifundefined{theorem}{\newtheorem{theorem}{Theorem}}{}
\@ifundefined{definition}{\newtheorem{definition}[theorem]{Definition}}{}
\@ifundefined{lemma}{\newtheorem{lemma}[theorem]{Lemma}}{}
\@ifundefined{proposition}{\newtheorem{proposition}[theorem]{Proposition}}{}
\makeatother
\newcommand{\crb}[1]{\left\{ #1 \right\}}
\newcommand{\di}{dimension}
\newcommand{\prs}[1]{\left( #1 \right)}
\title[Collapsibility in Multiparametric Models of Random Simplicia]{Collapsibility in Multiparametric Models of Random Simplicial Complexes}
\author{Jon V. Kogan}
\date{Source: arXiv:2606.15276v1 (CC BY 4.0)}
\begin{document}
\maketitle

\noindent\emph{Source and licence.} Collapsibility in Multiparametric Models of Random Simplicial Complexes. Version arXiv:2606.15276v1 (CC BY 4.0), distributed under the Creative Commons Attribution 4.0 International licence, as recorded in the arXiv licence metadata for this exact version and shown on the abstract page https://arxiv.org/abs/2606.15276v1. Full rights notice: licenses/arxiv-2606.15276-CC-BY-4.0.txt.\par\smallskip

\noindent\textit{Editorial scope.} Counted unit: the Lemma whose source label is \texttt{lem:upperBudgetCost} in the article (source0.tex lines 587--595, source label \texttt{lem:upperBudgetCost}), delivered together with the article's own complete original proof of it (source0.tex lines 596--628, one \texttt{proof} environment of 33 lines). No mathematics is written, repaired or re-derived: statement and proof are the article's own text line for line. Required context retained uncounted: FNNotation (lines 173--182); propBudCostLower (lines 458--461). Every definition, display and label of those lines is retained unchanged. Omitted from this package and not redistributed: 1--172, 183--457, 462--586, 629--748. Nothing inside a retained excerpt is omitted or altered: every retained body is delivered exactly as the article's lines are written, so no omission receipt of any kind accompanies this package. Citations: the retained excerpts carry no citation command at all, so no bibliography entry of the article is transcribed and no reference is redistributed. Reference adaptation: none was needed and none was made; every reference inside the delivered unit resolves inside it, so no unresolved reference or citation is produced. Printed slips kept literal: no printed word, symbol, hypothesis or number of the counted statement or of its proof was altered. Layout and class: the article's own class is \texttt{article} with packages including inputenc, authblk, setspace, geometry, graphicx, subcaption, amsmath, amssymb, amsthm, thmtools, tikz, tikz-cd, bbold, enumitem; the wrapper is the lane's pinned \texttt{amsart} editorial wrapper at 10pt on A4, which restates the theorem-like environments the retained text needs and adds the article's own macro definitions that the retained text needs (\texttt{\textbackslash crb}, \texttt{\textbackslash di}, \texttt{\textbackslash prs}), with extra wrapper packages (bm, bbm, dsfont, xspace, cancel, mathtools, cleveref, enumitem). The delivered numbering is the unit's own. Assets: no retained span contains a figure environment, a graphics-inclusion command, a PDF-page inclusion, a table or a TikZ picture; no image file is copied into this package, the package declares no asset dependency and no image file is redistributed with it. Licence: version arXiv:2606.15276v1 of this article is distributed under the Creative Commons Attribution 4.0 International licence, as recorded in the arXiv licence metadata for this exact version and shown on the abstract page https://arxiv.org/abs/2606.15276v1. A full rights notice is delivered at licenses/arxiv-2606.15276-CC-BY-4.0.txt. Characters: every retained excerpt is pure ASCII, so the delivered engine, XeLaTeX with its default Latin Modern text font, reports no missing character.
\begin{definition}\label{FNNotation}
    In $\overline{X}(n;p_1,p_2,...)$, where $p_i=n^{-\alpha_i}$:
    \begin{itemize}
        \item $\beta_i:=i+1-\alpha_i\sim \log_n E(\Delta^i\subset X)$, $\beta:=\max\{\beta_i\}$, $l=\lfloor\beta \rfloor$.
        \item $\gamma_k:=\max_{i\ge k}\{\beta_i\}\sim \log_n E(\Delta^k\subset \overline{X})$.
        \item $\nu_k:=2\gamma_k-k$,  $l':=\max\{k| \nu_k\ge 0\}\le \lfloor 2\beta \rfloor\le 2l+1$.
        \item $e_k:=\gamma_k-k=\nu_k-\gamma_k <0$ where the inequality is when $k>l$. $e_{k+1}\le e_k-1$.
    \end{itemize}
    (some facts are included for intuition's sake).
    \end{definition}

\begin{proposition}\label{propBudCostLower}
    The terms \textbf{budget} and \textbf{cost} are well defined for $\underline{X}(n;p_1,...)$ where $p_i=n^{-\alpha_i}$.
    More concretely, assume that a simplicial complex $A$ has $a_0$ vertices, $a_1$ edges and so on. Then $\lim_{n\to \infty}\log_n(E(A\subset \underline{X}))=a_0-\sum_{i=1} ^{\dim(A)}a_i\alpha_i$.
\end{proposition}

\begin{lemma}\label{lem:upperBudgetCost}
    Let $Z$ be a strongly connected complex w.r.t \di \ $l+1$. In $\overline{X}(n;p_1,...,p_r)$, $p_i=n^{
    -\alpha_i}$ we have the following:
    \begin{enumerate}
        \item $E(Z\subset \overline{X})=O(n^\beta)$. In particular, the budget is less than $l+1$.
        \item Let $Z'=Z\cup \tau, \tau\simeq\Delta^m, m>l$ be the result of an operation expanding $Z$ (w.r.t \di\ $l+1$) which increases the number of maximal faces. Then $E(Z'\subset_h \overline{X})=O(E(Z\subset_h \overline{X})n^{e_{l+1}+v+l-m})$, where $v$ is the number of new vertices added by the operation.
        In other words, the honest cost of such an operation is at least $m-e_{l+1}-v-l$ (note that $e_{l+1}<0$ and that $m\ge v+l$).
    \end{enumerate}
\end{lemma}

\begin{proof} \begin{enumerate}
    \item Consider a particular injective mapping $g$ of the vertices of $\sigma=\Delta^d, d\ge l+1$ to $[n]$. $g(\sigma)$ is a face in $\overline{X}$ if and only if $g(\sigma)$ is contained in a face in $X$ (the hypergraph).
Denote by $J_i$ the event that $g(\sigma)$ is contained in an $i$-\di al face in $X$. Then 
$$P(J_i)=\binom{n-d-1}{i-d}n^{-\alpha_i}-\sum_{j=0}^{i-d-1}\binom{n-d-1}{2(i-d)-j}n^{-2\alpha_i}+\dots=\theta(n^{i-\alpha_i-d})$$
where we use inclusion-exclusion, and the fact that the appearance of different simplexes in $X$ is independent. Since $\max(i+1-\alpha_i)=\beta<l+1$, we conclude, assuming $d\ge l+1$, that $i-\alpha_i-d<-1$, and so the first term dominates asymptotically.
Further, $c\sum_{j=0} ^\infty\frac{1}{n^j}=\theta(c)$, so the sum is of the claimed asymptotic class.
Using this we get 
$$P\prs{\bigcup_{i=d} ^r J_i}=\theta\prs{\sum_{i=d} ^r n^{i-\alpha_i-d}-\sum_{i\ge j\ge d} ^r n^{i+j-\alpha_i-\alpha_j-2d}+ \sum_{i\ge j\ge k\ge d}\dots}$$
where we use the same arguments again, and conclude 
$$P\prs{\bigcup_{i=d} ^r J_i}=\theta\prs{\sum_{i=d} ^r n^{i-\alpha_i-d}}.$$
By the same argument as in \Cref{propBudCostLower}, there are $\theta(n^{d+1})$ possible $g$'s, and so:
\begin{equation}\label{bugdetUpper}
    E(\Delta^{d}\subset \overline{X})=\theta(n^{d+1}(n^{-\alpha_d}+n^{1-\alpha_{d+1}}+...+n^{r-d-\alpha_r}))=\theta(n^{\max_{i\ge d}\{ i+1-\alpha_i\}})=\theta(n^{\gamma_d}).
\end{equation}

 But $\gamma_d\le \beta < l+1$, giving the budget.
 \item Denote $E(Z\subset_h\overline{X})=\zeta$, and suppose $Z'$ has $v$ vertices more than $Z$. Let $g$ be an injective map from the vertices of $Z'$ to $[n]$. $g(Z')$ is an honest subcomplex of $\overline{X}$ if and only if all the following are satisfied:
 \begin{enumerate}[label=\alph*)]
     \item $g(Z)$ is an honest subcomplex.
     \item $g(\tau)$ is a face in $\overline{X}$.
     \item No face of $\overline{X}$ contains both $g(\tau)$ and a maximal face of $g(Z)$.
 \end{enumerate}
 Let $\crb{Z_i}$ be the maximal faces of $Z$. Then by the same arguments as above
 $$P\prs{g(Z')\subset_h\overline{X}\mid g(Z)\subset_h\overline{X}}=$$$$\sum_{i=m} ^r \prs{\binom{n-m-1}{i-m}-\sum_{Z_j}\binom{n-|Z_j\cup \tau|}{i-|Z_j\cup \tau|+1}}n^{-\alpha_i}+o(n^{i-m-\alpha_i})=\theta(n^{\gamma_m-m-1}).$$
The difference of binomial coefficients is to exclude maximal faces containing both $g(\tau)$ and one of the $Z_j$ (containing more than one is excluded by the condition that $g(Z)$ is an honest subcomplex), and the asymptotic is achieved by the observation that $\sum_{Z_j}\binom{n-|Z_j\cup \tau|}{i-|Z_j\cup \tau|+1}=o\prs{\binom{n-m-1}{i-m}}$ as well.
$$P(g(Z')\subset_h\overline{X})=P(g(Z')\subset_h\overline{X}\mid g(Z)\subset_h\overline{X})\cdot P(g(Z)\subset_h\overline{X})=\theta\prs{P(g(Z)\subset_h\overline{X})n^{\gamma_m-m-1}}.$$
$E(Z\subset_h\overline{X})=\sum_{g|_Z}P(g(Z)\subset_h\overline{X})$ and similarly $E(Z'\subset_h\overline{X})=\sum_{g}P(g(Z')\subset_h\overline{X})$.
Therefore 
$$E(Z'\subset_h\overline{X})=\sum_{g}P(g(Z')\subset_h\overline{X})=\theta\prs{n^{\gamma_m-m-1}\sum_{g}P(g(Z)\subset_h\overline{X})}=$$
$$\theta\prs{n^{\gamma_m-m-1+v}\sum_{g|_Z}P(g(Z)\subset_h\overline{X})}=\theta\prs{n^{\gamma_m-m-1+v}E(Z\subset_h\overline{X})}.$$
Recall from \Cref{FNNotation} that $e_m=\gamma_m-m$, and that $e_m\le l+1-m+e_{l+1}$, and we get the desired result.
    \end{enumerate}
\end{proof}

\end{document}
